\section{Introduction}
\label{sec:intro}

The classical theory of piezoelectricity relates the electric displacement to
the strain and the electric field alone.  Flexoelectricity adds a dependence on
the \emph{gradient} of strain \cite{shen2010theory,deng2020impact}, so that a
bending deformation polarises a dielectric even when it is centrosymmetric and
carries no piezoelectric constant.  The effect is small in absolute terms but
it scales as the inverse of a length, and it therefore becomes competitive with
piezoelectricity in thin films, at free surfaces and around edges and cracks,
where the strain gradient is large.  In a semiconductor the polarisation is not
merely a passive consequence of deformation: the induced bound charge is
screened by the mobile carrier gas, the screening redistributes the carriers,
and the redistribution feeds back on the electric field.  This closed loop --
strain gradient $\rightarrow$ polarisation $\rightarrow$ carrier redistribution
$\rightarrow$ field $\rightarrow$ strain -- is the subject of the rapidly
developing theory of flexoelectric semiconductors
\cite{qu2020mobile,qu2022thermal}.

Two strands of that literature are directly relevant here.  The first is the
continuum formulation itself.  Starting from a free-energy density that contains
the strain, the strain gradient, the electric field and the carrier
concentration perturbation, one obtains the field equations by independent
variations with respect to the displacement, the potential and the carrier
density.  Qu, Jin and Yang \cite{qu2020mobile} used this route to show that the
mechanical field redistributes mobile charge in a composite flexoelectric
dielectric--semiconductor beam, Qu, Jin and Yang \cite{qu2022bending} carried the same construction to the
bending of a flexoelectric semiconductor plate, and Qu, Zhang, Gao and Jin
\cite{qu2022thermal} extended the model to thermally induced carrier
redistribution, where the temperature gradient acts as an additional driving
term for the charge.  The most complete statement of the framework is the
Coleman--Noll derivation of Qu, Pan, Zhu and Deng \cite{qu2025thermodynamic},
whose additive free-energy decomposition is the reference against which the
added term of \cref{sec:intro_coupling} is placed in
\cref{sec:disc_literature}.  All these works make the point that must be respected in any
subsequent extension: the extra constitutive couplings are not free parameters
of a model, they are constrained by a thermodynamic potential, and the
reciprocal partner of every added term in the stress must appear in the
electrochemical potential.

The second strand is the wave theory.  Yang, Deng, Liang and Shen
\cite{yang2018lamb} derived the dispersion of Lamb waves in a plate with both
flexoelectricity and strain-gradient elasticity, and Yang, Liang, Deng and Shen
\cite{yang2020rayleigh} treated Rayleigh waves on a centrosymmetric
flexoelectric half-space.  Wang, Zhang, Zhang, Du, Xiao and Lou
\cite{wang2025rayleigh} recently combined both effects with the carrier physics
of a piezoelectric semiconductor layer on an elastic substrate and showed that
the layer thickness, the steady carrier concentration and the bias field all
move the dispersion and the attenuation of the surface wave.  Biswas and Sahu \cite{biswas2025rayleigh} extended the
half-space treatment to an n-type flexoelectric piezoelectric semiconductor with
second-order strain gradients and reported how the flexoelectric coefficients and
the steady carrier concentration move the Rayleigh velocity of that structure.
Zhang, Zhao, Hu,
Nie and Liu \cite{zhang2025pscplate} gave an exact Stroh-type treatment of a
piezoelectric semiconductor plate and identified the range of carrier
concentrations in which the semiconductor coupling is active at all -- outside
that range the plate reverts to a piezoelectric or even a purely elastic
behaviour.  These studies establish three facts that frame the present work:
flexoelectric and strain-gradient terms change the dispersion of guided modes;
the carrier physics is active only inside a window of doping; and the
attenuation of a guided mode can be controlled by the electrical boundary
conditions rather than by the mechanical ones.

What has not been examined in that literature is the \emph{scattering} problem
at an edge of a guided plate.  Reflection at a free surface of a flexoelectric
medium has been treated for plane waves in a half-space, where the result is a
set of reflection coefficients for the few body waves that the half-space
supports \cite{mondal2022scattering}; the modal partition of a \emph{guided}
plate branch at an edge, where the outgoing spectrum is discrete and contains a
backward-propagating member, and where the carriers participate in the
partition, has not been computed.  A guided mode that travels towards a free
edge has to convert into the set of outgoing modes that the edge can support.  In a purely elastic plate
this is a classical problem: an incident mode is reflected into its time-reverse
partner and into a finite family of mode-converted branches, and a
backward-propagating branch --- one whose energy flux opposes its phase
velocity --- can be excited whenever the edge problem admits such a branch.  We
therefore ask two questions and answer them separately.  Does this plate admit a
backward branch at the reported frequency, and how strongly does the edge
convert into it?  And does the carrier--strain-gradient coupling bias that
partition?  The answers turn out to be different in kind: the branch, and
therefore the strong conversion, is supplied by the strain-gradient sector
alone, while the added coupling only modifies it --- by one per cent at the
operating point, and by a third at the certified maximum through the extra
attenuation it gives the backward mode.  Reporting the two separately, rather
than attributing the effect to the new term, is what the rest of the paper
does.

\subsection{The added coupling and its reciprocals}
\label{sec:intro_coupling}

The model used here is deliberately minimal.  On top of the standard
flexoelectric piezoelectric-semiconductor enthalpy we add a single bilinear
term,
\begin{equation}
\label{eq:added_coupling}
\Delta\mathcal{H}=\Delta n\,H_{jkl}\,\varepsilon_{jk,l},
\end{equation}
where $\Delta n$ is the carrier-concentration perturbation, $\varepsilon_{jk,l}$
the strain gradient and $H_{jkl}$ a third-order coupling tensor with the units
of energy times length.  Because the enthalpy is a thermodynamic potential, the
term in \cref{eq:added_coupling} generates two partners.  Differentiating with
respect to the strain gradient gives an addition to the higher-order
(hyper)stress,
\begin{equation}
\label{eq:added_hyperstress}
\Delta\tau_{jkl}=\frac{\partial\Delta\mathcal{H}}{\partial\varepsilon_{jk,l}}
=H_{jkl}\,\Delta n,
\end{equation}
and differentiating with respect to the carrier density gives an addition to the
electrochemical potential,
\begin{equation}
\label{eq:added_mu}
\Delta\widetilde\mu=\frac{\partial\Delta\mathcal{H}}{\partial\Delta n}
=H_{jkl}\,\varepsilon_{jk,l}.
\end{equation}
Equations \eqref{eq:added_coupling}--\eqref{eq:added_mu} are one assumption,
not three.  They are the only new physics in the paper.  No additional field,
no micro-inertia, no flexoelectric-gradient term and no nonlinear
electromechanical feedback is introduced, and the classical constants are not
refitted.

The reciprocal relation \eqref{eq:added_mu} is what makes the edge problem
interesting.  A guided plate mode carries a strain gradient through the
thickness, so it perturbs the electrochemical potential even before any carrier
moves; the carriers respond, the charge density changes, and the changed charge
density modifies the hyperstress through \eqref{eq:added_hyperstress}.  The loop
is closed at the edge, where the conversion into outgoing branches takes place.

\subsection{Scope and organisation}
\label{sec:intro_scope}

The paper is organised as follows.  \Cref{sec:governing} states the
variational principle, the constitutive relations, the field equations and the
boundary conditions, and fixes the nomenclature for the plate surfaces, the
reflecting edge, the mechanical tractions, the higher-order double tractions,
the electrical conditions and the carrier-blocking conditions.
\Cref{sec:numerics} describes the solution strategy: a Fourier reduction in the
tangential direction with a single conserved wavenumber $k_y$, a spectral
Galerkin discretisation through the thickness, a companion-matrix formulation
in the direction of propagation, and an edge-matching step that separates
incoming from outgoing modes.
\Cref{sec:results} presents the ten figures of the study and, for each, states
what the panel is allowed to claim.  \Cref{sec:discussion} discusses what the
numbers mean and, at least as importantly, what they do not mean.
\Cref{sec:conclusions} lists the limitations that we could not remove.

Three conventions are used throughout and are worth stating once.  The plate
occupies $x<0$ and the reflecting edge is at $x=0$; the $z$ axis is the
six-fold axis of ZnO and points through the thickness, with the faces at
$z=\pm d/2$.  A quantity is called \emph{backward} when its energy flux has a
negative component along $x$, which is the direction of propagation of the
incident wave.  A mode is called \emph{outgoing} when it either decays away
from the edge or, if it is propagating, carries energy away from the edge.
All three definitions are operational: they are the ones implemented in the
solver, not merely the ones used in the prose.
